\clearpage
\VerdierPhysicalPageAnchor{0161}%
\begingroup\small\setlength{\parskip}{0pt}\setlength{\emergencystretch}{1.5em}
\VerdierSemanticAnchor{ch3-3-2-9-exactness-test}
est exacte dans $\mathcal D^\wedge$. La suite (3.2.9.1) est exacte si et
seulement si, pour tout objet $W$ de $\mathcal D$, la suite
\[
\operatorname{Hom}_{\mathcal D^\wedge}(h(W),h(X))\longrightarrow
\operatorname{Hom}_{\mathcal D^\wedge}(h(W),h(Y))\longrightarrow
\operatorname{Hom}_{\mathcal D^\wedge}(h(W),h(Z))
\]
est exacte. Or ceci résulte de ce que $h$ est pleinement fidèle et de (1.2.1).

\VerdierSemanticAnchor{ch3-3-2-9-factorization}
Soit alors $\mathcal A$ une catégorie abélienne et
\[
\mathcal E x(\mathrm A(\mathcal D),\mathcal A)
\longrightarrow\mathcal H\!om(\mathcal D,\mathcal A)
\]
le foncteur « composition avec $\mathcal H$ ». Il est clair que ce foncteur
est pleinement fidèle : un morphisme entre foncteurs exacts est uniquement
déterminé lorsqu'on connaît ses valeurs sur suffisamment d'objets
projectifs. Il reste à démontrer que tout foncteur cohomologique
$\Psi\colon\mathcal D\to\mathcal A$ se factorise à travers
$\mathrm A(\mathcal D)$ :
\VerdierSemanticAnchor{ch3-3-2-9-2}
\[
\begin{tikzcd}[column sep=5em,row sep=3em]
\mathcal D\arrow[r,"\mathcal H"]\arrow[dr,"\Psi"']&
\mathrm A(\mathcal D)\arrow[d,"\Phi"]\\
&\mathcal A
\end{tikzcd}
\tag{3.2.9.2}
\]
le foncteur $\Phi$ étant exact. Pour tout objet
$X\xrightarrow fY$ de $\mathrm A(\mathcal D)$, posons
\begingroup\renewcommand{\thefootnote}{(*)}%
\[
\Phi(X\xrightarrow fY)=
\operatorname{Image}_{\mathcal A}
(\Psi(X)\xrightarrow{\Psi(f)}\Psi(Y)).\footnote{On remarquera que la
construction du foncteur $\Phi$ fait intervenir le choix d'un foncteur
« image » dans $\mathcal A$. Ce choix est bien entendu unique à isomorphisme
unique près. On peut toujours imposer, ce que nous ferons désormais, que le
foncteur « image » associe à tout morphisme identique
$\operatorname{id}_T$ dans $\mathcal A$ l'objet $T$ muni du morphisme
$T\xrightarrow{\operatorname{id}_T}T$.}
\]
\endgroup
On détermine ainsi un foncteur de $\mathrm A(\mathcal D)$ dans $\mathcal A$
tel que le diagramme (3.2.9.2) soit commutatif. Montrons que $\Phi$ est un
foncteur exact. Il résulte immédiatement de (3.2.7) que $\Phi$ est exact à
gauche. En passant aux catégories opposées, on montre de même que $\Phi$ est
exact à droite. Ceci achève la démonstration du théorème (3.2.1).

\VerdierSemanticAnchor{ch3-3-2-10}
\medskip\noindent\textbf{Remarque 3.2.10.}\quad
1) La propriété universelle décrite dans le théorème (3.2.1, c)) détermine
la catégorie $\mathrm A(\mathcal D)$ à équivalence près.

2) Lorsque dans $\mathcal D$ les sommes dénombrables ou bien les produits
dénombrables sont représentables, tous les objets projectifs
(resp. injectifs) de $\mathrm A(\mathcal D)$ sont isomorphes à des objets de
$\mathcal D$. En effet, un tel objet est facteur direct dans
$\mathrm A(\mathcal D)$ d'un objet de $\mathcal D$. Il est donc noyau d'un
projecteur d'un objet de $\mathcal D$. Or, sous les hypothèses faites, les
projecteurs sont décomposables dans $\mathcal D$ (1.2.9).
\endgroup
